\paragraph{Complete-pipeline comparison.}\AchievementDetectionEN{} Paired prompt-group uncertainty is reported in Table~\ref{tab:uncertainty}.

\begin{table}[t]
\centering\footnotesize
\caption{In-domain detection (\%). Mean $\pm$ sample SD across available seeds; a single completed seed has no SD. Thresholds use only in-domain calibration.}\label{tab:main}
\begin{tabular}{lrrrr}\toprule
Method & R@128 & FPR & AUROC & AP \\
\midrule
Prompt LR & 22.9 $\pm$ 0.0 & 3.3 $\pm$ 0.0 & 74.4 $\pm$ 0.0 & 80.1 $\pm$ 0.0 \\
Prefix TF--IDF LR & 28.2 $\pm$ 0.0 & 1.6 $\pm$ 0.0 & 81.1 $\pm$ 0.0 & 86.3 $\pm$ 0.0 \\
Current-state LR & 43.7 $\pm$ 0.0 & 5.6 $\pm$ 0.0 & 83.1 $\pm$ 0.0 & 86.5 $\pm$ 0.0 \\
Running-mean LR & 46.1 $\pm$ 0.0 & 6.2 $\pm$ 0.0 & 83.5 $\pm$ 0.0 & 86.5 $\pm$ 0.0 \\
Dynamic GRU & 38.1 $\pm$ 6.3 & 3.3 $\pm$ 1.0 & 83.6 $\pm$ 0.9 & 87.6 $\pm$ 0.8 \\
\bottomrule\end{tabular}
\end{table}


\paragraph{Detector trade-offs.}Table~\ref{tab:main} reports the full baseline comparison. Static probes, MLPs and recurrent heads offer complementary recall--FPR profiles under the common calibration target, supporting application-specific detector selection.

\begin{table}[t]
\centering\footnotesize
\caption{In-domain alert positions. Timing is conditional on true detections; recall retains all unsafe responses in its denominator.}\label{tab:alert_positions}
\begin{tabular}{lrrrr}\toprule
Method & Recall & Mean $\tau$ & Med. $\tau$ & $\tau/T$ (\%) \\
\midrule
Prompt LR & 22.9 $\pm$ 0.0 & 8.0 $\pm$ 0.0 & 8.0 $\pm$ 0.0 & 14.8 $\pm$ 0.0 \\
Prefix TF--IDF LR & 28.3 $\pm$ 0.0 & 20.7 $\pm$ 0.0 & 8.0 $\pm$ 0.0 & 26.2 $\pm$ 0.0 \\
Current-state LR & 43.9 $\pm$ 0.0 & 27.5 $\pm$ 0.0 & 16.0 $\pm$ 0.0 & 34.6 $\pm$ 0.0 \\
Running-mean LR & 46.3 $\pm$ 0.0 & 15.1 $\pm$ 0.0 & 8.0 $\pm$ 0.0 & 23.5 $\pm$ 0.0 \\
Dynamic GRU & 38.4 $\pm$ 6.1 & 47.9 $\pm$ 15.1 & 41.8 $\pm$ 18.3 & 54.0 $\pm$ 11.5 \\
\bottomrule\end{tabular}
\end{table}


\begin{table}[t]
\centering\footnotesize
\caption{Prompt-group paired bootstrap: Recall@128 differences and 95\% intervals (percentage points).}\label{tab:uncertainty}
\begin{tabular}{lp{.34\columnwidth}rr}\toprule
Split & Contrast & $\Delta$ & 95\% CI \\
\midrule
test & Dynamic GRU vs. Prompt LR & +15.2 & [+12.7, +17.8] \\
test & Dynamic GRU vs. Current-state LR & -5.6 & [-7.1, -3.9] \\
test & Dynamic GRU vs. Prefix TF--IDF LR & +9.9 & [+7.6, +12.2] \\
external & Dynamic GRU vs. Prompt LR & +12.3 & [+6.7, +18.8] \\
external & Dynamic GRU vs. Current-state LR & -11.7 & [-20.6, -3.6] \\
external & Dynamic GRU vs. Prefix TF--IDF LR & +12.3 & [+6.7, +18.8] \\
\bottomrule\end{tabular}
\end{table}


\paragraph{Same-prompt outcomes.}The test subset containing both final labels per prompt has 97 groups and 207 responses. Dynamic GRU and current-state LR attain 20.1\% and 22.3\% Recall@128, at 6.1\% and 6.7\% FPR, respectively, using unchanged thresholds. The matched-prompt subset characterizes how the two detectors respond to different outcomes associated with a shared question.

\paragraph{History-aggregation control.}Running-mean LR obtains 46.1\% Recall@128; the dynamic-GRU difference is -8.1 percentage points. The comparison complements current-state and learned-recurrence baselines by quantifying causal history accumulation.

\paragraph{Uncertainty.}Table~\ref{tab:uncertainty} retains paired prompt-group bootstrap intervals. Each paired draw preserves the prompt groups and fixed fitted seeds across the compared detectors.

\paragraph{Pre-specified contrast.}The dynamic-GRU minus static-probe Recall@128 difference is -5.64 percentage points (95\% prompt-group interval [-7.10, -3.93]). The static probe favors recall coverage, while the GRU has the lower observed test FPR; these operating profiles support different monitoring priorities.
